% ═════════════════════════════════════════════════════════════════════════════
%  第五章  讨论
% ═════════════════════════════════════════════════════════════════════════════
\section{讨论}
\label{sec:discussion}

% ─────────────────────────────────────────────────────────────────────────────
\subsection{各检索路径贡献分析}
\label{sec:disc-path}

三路并行检索架构（Dense / BM25 / KG）的核心假设是：
不同路径分别覆盖装配问答中的不同信息需求。
基于第~\ref{sec:exp-main}~节与~\ref{sec:exp-ablation}~节的实验数据，
本文将其作用概括为"语义召回、精确匹配、结构约束"三类互补能力。

\textbf{Dense 与 BM25 的互补性。}~密集向量检索（B1 Dense-only）擅长捕捉
"压气机前支撑轴"这类语义近邻查询，
在 Tier-1 精评的 Hit@1 上达到 0.983，体现出较强的首位命中能力。
BM25（B2）则在含精确零件号（如 \texttt{8210-210A}）、
扭矩数值（\texttt{32-36\,lb.in}）的字符串匹配场景上具有天然优势。
RRF 融合（B3 Hybrid）保留二者优势，
使系统既能处理同义改写，也能保留工业文档中不可模糊化的编号与数值。

\textbf{KG 路径的结构化推理。}~B4 MHRAG 在 Tier-1 生成阶段取得
Relevance 3.50、Completeness 3.02 与 Correctness 2.35。
其中 Correctness 相对 B3 Hybrid 提升 14.6\%，
说明 KG 子图对事实核验和关系补全有直接帮助。
进一步观察题型分组，KG 路径在工序工具类与技术规范类问题上贡献最明显：
前者依赖 \textit{requires}、\textit{precedes} 等工序关系，
后者依赖 \textit{specifiedBy} 对扭矩、间隙和材料要求的结构化聚合。
这也是本文三源 KG 相对纯文本检索的关键价值。

\textbf{跨章节场景的价值。}~在 Tier-2 跨章节粗评中，
B4 MHRAG 在 Relevance、Completeness 与 Correctness 三项指标上均取得数值最高值。
其核心原因并不是 KG 替代了文本检索，
而是 KG 将 BOM 中的零件身份、CAD/IPC 中的层级关系和手册中的工序知识统一到同一实体空间，
从而缓解了单一章节关键词检索容易漏掉关系链的问题。

\textbf{Reranker 的定位。}~本文保留 BGE-Reranker-Base 作为粗排后的精排模块，
用于在 Dense、BM25 与 KG 子图之间重新选择最终上下文。
当前实验未单独量化移除 Reranker 后的三路配置表现，
因此本文不将 Reranker 的收益作为主要创新结论；
它更适合作为工程系统中的质量控制组件。

% ─────────────────────────────────────────────────────────────────────────────
\subsection{知识图谱质量对生成效果的影响}
\label{sec:disc-kg-quality}

\subsubsection{三源 KG 构建效果}

本文构建的 PT6A 整机知识图谱共 7{,}333 节点 / 8{,}752 边，
9 类核心语义关系 + 2 类辅助关系（SAME\_AS / interchangesWith）均有覆盖。
其中：
\textit{isPartOf} 2{,}869 条（BOM + CAD），
\textit{participatesIn} 1{,}750 条（手册），
\textit{precedes} 1{,}241 条（手册），
\textit{adjacentTo} 931 条（CAD 装配兄弟邻接），
\textit{specifiedBy} 649 条、\textit{matesWith} 559 条等。
本体设计涵盖完整的 7 类实体节点
（Part / Assembly / Procedure / Tool / Specification / Interface / GeoConstraint），
表明三源融合 KG 在覆盖度上相对单一文本来源的工作具有更全面的语义表达。

\subsubsection{跨源实体对齐挑战}

PT6A 维修手册原文为英文，而黄金三元组集采用中文实体名（专业术语翻译），
两者在系统检索阶段需经过 BGE-M3 多语言对齐桥接。
实验中观察到部分专业术语（如"细手锉刀"\,$\rightarrow$\,Hand File Fine）
的语言对应在 BGE-M3 默认设置下未能稳定召回，
为此本文在 QA 评测时引入双语 query 联合检索策略
（\texttt{question + question\_en} 拼接），有效缓解了此问题。
这提示工业 RAG 系统在多语种场景下应内置查询自动翻译与领域术语词典。

\subsubsection{HITL 审阅的边界价值}

第~\ref{sec:hitl}~节的 HITL 审阅机制在本文实验中并未量化测试。
其设计动机在于：
对低置信度三元组与跨源命名歧义条目的领域专家校验，
是无法被纯自动化流水线替代的。
机制本身作为 KG 构建管道的可控接口提供端到端工程价值，
其量化收益依赖于具体团队的领域专家可用度，
不在本文实验主表的核心指标范围内。

% ─────────────────────────────────────────────────────────────────────────────
\subsection{局限性与未来工作}
\label{sec:disc-limit}

\textbf{评测集规模与来源。}~Tier-1 精评集来自单一压气机子系统，
且问题由黄金三元组模板反向生成，适合验证结构化知识覆盖，
但仍不能充分代表真实生产环境中的自然提问分布。
未来应补充领域专家自由提问、多人独立标注和生产现场问答日志，
以降低模板化评测带来的偏置。

\textbf{Judge 评分的外部效度。}~Tier-2 跨章节粗评（20 题）依赖 LLM-Judge 单一指标，
不含 \texttt{gold\_chunk} 标注。
该评测方式适合验证系统覆盖广度，
但在具体题目上的判别精度受 Judge 模型自身能力影响。
未来工作可邀请领域专家对 Tier-2 答案进行二次盲审，
建立专家-Judge 一致性矩阵以校准 Judge 评分偏置。

\textbf{语料语言局限。}~本文实验语料 PT6A 维修手册为英文原文，
评测黄金集为中文专业术语翻译。
中英对应在 BGE-M3 的多语言能力下大体可工作，
但存在专业术语词典稀缺导致的部分召回失败案例。
未来工作可引入领域术语对齐词典或在检索阶段调用 LLM 进行查询改写。

\textbf{CAD 几何邻接的精度。}~由于实验环境未启用 \texttt{pythonocc-core}，
\textit{adjacentTo} 关系采用降级实现（同父装配下的兄弟节点视为逻辑邻接），
而非论文§\ref{sec:kg}~描述的轴对齐包围盒（AABB）间距计算。
该简化在装配树清晰的 PT6A IPC 文档下保持可用性，
但在复杂多层装配场景下可能产生误判。
未来工作将在生产环境集成 OCC 后回归 AABB 物理邻接计算。
